% Subsection 5.2 of §5 Leakage. Begins at \subsubsection level under \subsection{International reallocation}.

The international margin asks whether Belgian buyers shift purchases of regulated products away from EU sources (where production is subject to the EU ETS) toward non-EU sources (where it is not). This is the canonical leakage channel that motivates CBAM. On the customs panel 2000--2022, we measure the regulated and unregulated import shares from EU and non-EU sources at the HS6-product level.

\paragraph{Setup.} For each HS6 product $p$, we classify the product as \emph{regulated} if it maps to a NACE 4-digit category that contains at least one ETS installation, and \emph{unregulated} otherwise. We separately track imports from EU-ETS source countries and from non-ETS source countries. The leakage hypothesis predicts that, post-MSR, Belgian regulated-product imports should shift from EU-ETS sources toward non-ETS sources. We replicate the \citet{CosterMejeanDiGiovanni2024} share specification on Belgian customs over 2000--2022, extending their 2000--2019 window through the 2020--2022 EUA peak.

\paragraph{Headline result.} Figure~\ref{fig:cdgm_figure2} reports the four within-product-type share trajectories. The headline fact is in the top-left panel: the share of regulated imports sourced from ETS countries is essentially flat over the entire 2000--2022 window, including the post-MSR period in which the EUA price rose from \EUR{7} to peaks above \EUR{90}. If the EU ETS were inducing Belgian buyers to substitute away from EU sources of regulated products, this is the panel in which it would appear. It does not.

\begin{figure}[H]
    \centering
    \includegraphics[width=0.95\textwidth]{phase2_cdgm_figure2.png}
    \caption{Belgian import shares by product regulation status and source country regulation status, 2000--2022. The regulated $\times$ ETS share (top-left) is the panel in which carbon leakage to non-EU sources would appear; it is flat across the post-MSR run-up. The visible cross-bloc substitution is in the unregulated $\times$ ETS share (top-right, declining) and unregulated $\times$ non-ETS share (bottom-right, rising) --- which is wrong-signed for a carbon-policy mechanism. Source: NBB customs panel; product regulation status from EU ETS coverage rules.}
    \label{fig:cdgm_figure2}
\end{figure}

\paragraph{Reading the figure.} The two panels that do move are the unregulated ones: the share of unregulated imports from EU sources falls, and the share from non-EU sources rises. The pattern is the opposite of what carbon leakage would predict --- where leakage operates on the product margin most exposed to the EU ETS, here it is the products outside the EU ETS that are reallocating away from EU origins. The natural explanation is that the visible cross-bloc shift in unregulated imports reflects generic globalization of unregulated and commoditized product categories --- consistent with the China shock literature on Belgian manufacturing and with the rise of low-cost non-EU intermediates over 2000--2022 --- and is unrelated to carbon policy.

\paragraph{Belgium vs.\ France.} The Belgian flat share contrasts with \citet{CosterMejeanDiGiovanni2024}'s French finding of a Phase~3 share coefficient of $+0.121$ (significant at $p < 0.001$) on the same specification. The Belgian Phase~3 coefficient on the corresponding regression is approximately two orders of magnitude smaller and points in the opposite direction. The result is robust across alternative control-group definitions and fixed-effect structures. A national leakage outcome from a common EU policy regime is, therefore, not a universal feature of EU member states --- it is contingent on the importing economy's industrial structure and trade integration. Three plausible drivers of the Belgium--France divergence are intra-EU integration density (Belgian buyers consolidate to EU-domestic regulated suppliers rather than pivot to non-ETS countries), pre-existing chemicals saturation (the Antwerp cluster was heavily integrated with non-ETS exporters before 2005, leaving limited headroom), and country-level differences in sectoral composition. Distinguishing them requires variation across small open EU economies and is left for future work.

\paragraph{Reading the international and domestic margins together.} Neither the within-Belgium reallocation tests of \S\ref{sec:domestic} nor the cross-border test here find a leakage signal. The within-Belgium tests rule out substitution among regulated suppliers and across input categories; the cross-border test rules out substitution toward non-EU sources of regulated products. The set of reallocation margins on which we can find leakage in Belgium under the post-MSR EU ETS is empty. Why is the subject of the next section.
